\section{Silver's theorem}

\import{silver/last-lemma.tex}
\import{cardinal/normal-sequence.tex}
\import{set/symdiff.tex}

The final coding sends a subset of the singular cardinal to its intersections
with a normal cofinal sequence.  Each graph below is a set obtained by bounded
Replacement; no class-valued cardinal enumeration is used.

\begin{definition}\label{silver_powerset_fibres}
    $\silverPowersetFibres{c}{\lambda}
        =\{(i,\pow{c(i)})\mid i\in\lambda\}$.
\end{definition}

\begin{definition}\label{silver_intersection_code}
    $\silverIntersectionCode{c}{X}{\lambda}
        =\{(i,X\inter c(i))\mid i\in\lambda\}$.
\end{definition}

\begin{definition}\label{silver_intersection_map}
    $\silverIntersectionMap{c}{\kappa}{\lambda}
        =\{(X,\silverIntersectionCode{c}{X}{\lambda})
            \mid X\in\pow{\kappa}\}$.
\end{definition}

\begin{proposition}\label{silver_powerset_fibres_type}
    Let $A=\silverPowersetFibres{c}{\lambda}$.
    Then $A$ is a function on $\lambda$ and
        for every $i\in\lambda$ we have $A(i)=\pow{c(i)}$.
\end{proposition}
\begin{proof}
    $A$ is a relation by \auto{relation,silver_powerset_fibres}.
    $A$ is right-unique
        by \auto{rightunique,silver_powerset_fibres,pair_eq_iff}.
    $\dom{A}=\lambda$
        by \auto{setext,dom_iff,silver_powerset_fibres,pair_eq_iff}.
    Hence $A$ is a function on $\lambda$.
    For every $i\in\lambda$ we have $(i,\pow{c(i)})\in A$
        by \obv{silver_powerset_fibres}.
    Hence for every $i\in\lambda$ we have $A(i)=\pow{c(i)}$
        by \obv{function_apply_intro}.
\end{proof}

\begin{proposition}\label{silver_intersection_code_type}
    Suppose $c$ is a function on $\lambda$.
    Let $A=\silverPowersetFibres{c}{\lambda}$.
    Then $\silverIntersectionCode{c}{X}{\lambda}
        \in\dependentProduct{\lambda}{A}$ and
        for every $i\in\lambda$ we have
            $(\silverIntersectionCode{c}{X}{\lambda})(i)=X\inter c(i)$.
\end{proposition}
\begin{proof}
    Let $h=\silverIntersectionCode{c}{X}{\lambda}$.
    $A$ is a function on $\lambda$ and
        for every $i\in\lambda$ we have $A(i)=\pow{c(i)}$
        by \obv{silver_powerset_fibres_type}.
    $h$ is a relation by \auto{relation,silver_intersection_code}.
    $h$ is right-unique
        by \auto{rightunique,silver_intersection_code,pair_eq_iff}.
    $\dom{h}=\lambda$
        by \auto{setext,dom_iff,silver_intersection_code,pair_eq_iff}.
    For every $i\in\lambda$ we have $(i,X\inter c(i))\in h$
        by \obv{silver_intersection_code}.
    Hence for every $i\in\lambda$ we have $h(i)=X\inter c(i)$
        by \obv{function_apply_intro}.
    For every $i\in\lambda$ we have $X\inter c(i)\subseteq c(i)$
        by \obv{inter_lower_right}.
    Hence for every $i\in\lambda$ we have $h(i)\in A(i)$
        by \auto{powerset_intro}.
    Hence for every $i\in\lambda$ we have
        $h(i)\in\unions{\ran{A}}$
        by \obv{indexed_family_member_union}.
    Hence $h$ is a function to $\unions{\ran{A}}$.
    Hence $h\in\funs{\lambda}{\unions{\ran{A}}}$
        by \obv{funs_intro}.
    Hence $h\in\dependentProduct{\lambda}{A}$
        by \obv{dependent_product_iff}.
    Follows by \auto.
\end{proof}

\begin{proposition}\label{silver_intersection_code_application}
    Suppose $c$ is a function on $\lambda$.
    Let $i\in\lambda$.
    Then $(\silverIntersectionCode{c}{X}{\lambda})(i)=X\inter c(i)$.
\end{proposition}
\begin{proof}
    Let $A=\silverPowersetFibres{c}{\lambda}$.
    Follows by \auto{silver_intersection_code_type}.
\end{proof}

\begin{proposition}\label{silver_intersection_map_function}
    Suppose $c$ is a function on $\lambda$.
    Let $A=\silverPowersetFibres{c}{\lambda}$.
    Let $e=\silverIntersectionMap{c}{\kappa}{\lambda}$.
    Then $e\in\funs{\pow{\kappa}}{\dependentProduct{\lambda}{A}}$ and
        for every $X\in\pow{\kappa}$ we have
            $e(X)=\silverIntersectionCode{c}{X}{\lambda}$.
\end{proposition}
\begin{proof}
    $e$ is a relation by \auto{relation,silver_intersection_map}.
    $e$ is right-unique
        by \auto{rightunique,silver_intersection_map,pair_eq_iff}.
    $\dom{e}=\pow{\kappa}$
        by \auto{setext,dom_iff,silver_intersection_map,pair_eq_iff}.
    For every $X\in\pow{\kappa}$ we have
        $(X,\silverIntersectionCode{c}{X}{\lambda})\in e$
        by \obv{silver_intersection_map}.
    Hence for every $X\in\pow{\kappa}$ we have
        $e(X)=\silverIntersectionCode{c}{X}{\lambda}$
        by \obv{function_apply_intro}.
    For every $X\in\pow{\kappa}$ we have
        $\silverIntersectionCode{c}{X}{\lambda}
            \in\dependentProduct{\lambda}{A}$
        by \auto{silver_intersection_code_type}.
    Hence $e$ is a function to $\dependentProduct{\lambda}{A}$
        by \auto.
    Hence $e\in\funs{\pow{\kappa}}{\dependentProduct{\lambda}{A}}$
        by \obv{funs_intro}.
\end{proof}

\begin{proposition}\label{distinct_subsets_have_symmetric_difference}
    Suppose $X\subseteq\kappa$ and $Y\subseteq\kappa$.
    Suppose $X\neq Y$.
    Then $X\symdiff Y$ is inhabited and $X\symdiff Y\subseteq\kappa$.
\end{proposition}
\begin{proof}
    Show that $X\symdiff Y$ is inhabited.
    \begin{subproof}
        Suppose not.
        Show that every element of $X$ is an element of $Y$.
        \begin{subproof}
            Fix $a\in X$.
            Show that $a\in Y$.
            \begin{subproof}
                Suppose not.
                Hence $a\in X\setminus Y$ by \obv{setminus}.
                Hence $a\in X\symdiff Y$ by \auto{symdiff,union_iff}.
                Contradiction by \auto.
            \end{subproof}
        \end{subproof}
        Hence $X\subseteq Y$ by \obv{subseteq}.
        Show that every element of $Y$ is an element of $X$.
        \begin{subproof}
            Fix $a\in Y$.
            Show that $a\in X$.
            \begin{subproof}
                Suppose not.
                Hence $a\in Y\setminus X$ by \obv{setminus}.
                Hence $a\in X\symdiff Y$ by \auto{symdiff,union_iff}.
                Contradiction by \auto.
            \end{subproof}
        \end{subproof}
        Hence $Y\subseteq X$ by \obv{subseteq}.
        Hence $X=Y$ by \obv{subseteq_antisymmetric}.
        Contradiction.
    \end{subproof}
    $X\setminus Y\subseteq X$ and $Y\setminus X\subseteq Y$
        by \obv{setminus_subseteq}.
    Hence $X\setminus Y\subseteq\kappa$ and
        $Y\setminus X\subseteq\kappa$
        by \obv{subseteq_transitive}.
    Hence $(X\setminus Y)\union(Y\setminus X)\subseteq\kappa$
        by \auto{subseteq_union_iff}.
    Hence $X\symdiff Y\subseteq\kappa$ by \obv{symdiff}.
\end{proof}

\begin{proposition}\label{silver_intersection_codes_differ_at}
    Suppose $c$ is a function on $\lambda$.
    Let $i\in\lambda$.
    Suppose $a\in X\symdiff Y$ and $a\in c(i)$.
    Then $(\silverIntersectionCode{c}{X}{\lambda})(i)
        \neq(\silverIntersectionCode{c}{Y}{\lambda})(i)$.
\end{proposition}
\begin{proof}
    $(\silverIntersectionCode{c}{X}{\lambda})(i)=X\inter c(i)$ and
        $(\silverIntersectionCode{c}{Y}{\lambda})(i)=Y\inter c(i)$
        by \obv{silver_intersection_code_application}.
    $a\in X\setminus Y$ or $a\in Y\setminus X$
        by \obv{symdiff,union_iff}.
    Show that $X\inter c(i)\neq Y\inter c(i)$.
    \begin{subproof}
        Suppose $X\inter c(i)=Y\inter c(i)$.
        \begin{byCase}
            \caseOf{$a\in X\setminus Y$.}
                Hence $a\in X$ and $a\notin Y$ by \obv{setminus}.
                Hence $a\in X\inter c(i)$ by \obv{inter_intro}.
                Hence $a\in Y\inter c(i)$.
                Hence $a\in Y$ by \obv{inter_elim_left}.
                Contradiction.
            \caseOf{$a\notin X\setminus Y$.}
                Hence $a\in Y\setminus X$.
                Hence $a\in Y$ and $a\notin X$ by \obv{setminus}.
                Hence $a\in Y\inter c(i)$ by \obv{inter_intro}.
                Hence $a\in X\inter c(i)$.
                Hence $a\in X$ by \obv{inter_elim_left}.
                Contradiction.
        \end{byCase}
    \end{subproof}
\end{proof}

\begin{theorem}\label{silver_intersection_map_injection}
    Let $\kappa$ be a limit ordinal.
    Suppose $\cofinalMap{c}{\lambda}{\kappa}$.
    Let $A=\silverPowersetFibres{c}{\lambda}$.
    Let $e=\silverIntersectionMap{c}{\kappa}{\lambda}$.
    Then $e\in\inj{\pow{\kappa}}{\dependentProduct{\lambda}{A}}$.
\end{theorem}
\begin{proof}
    $c\in\funs{\lambda}{\kappa}$ by \obv{cofinal_map}.
    Hence $c$ is a function and $\dom{c}=\lambda$
        by \obv{funs_is_function,funs}.
    Hence $c$ is a function on $\lambda$.
    $e\in\funs{\pow{\kappa}}{\dependentProduct{\lambda}{A}}$
        by \auto{silver_intersection_map_function}.
    $\ordinalUnbounded{\ran{c}}{\kappa}$
        by \obv{cofinal_map_into_limit_has_unbounded_range}.
    Show that for all $X,Y\in\pow{\kappa}$ such that $X\neq Y$ we have
        $e(X)\neq e(Y)$.
    \begin{subproof}
        Fix $X,Y\in\pow{\kappa}$ such that $X\neq Y$.
        Hence $X\subseteq\kappa$ and $Y\subseteq\kappa$
            by \obv{powerset_elim}.
        Hence $X\symdiff Y$ is inhabited and
            $X\symdiff Y\subseteq\kappa$
            by \obv{distinct_subsets_have_symmetric_difference}.
        Take $a$ such that $a\in X\symdiff Y$.
        Hence $a\in\kappa$ by \obv{subseteq}.
        Take $y\in\ran{c}$ such that $a\in y$
            by \obv{ordinal_unbounded}.
        Take $i\in\dom{c}$ such that $c(i)=y$
            by \obv{fun_ran_iff}.
        Hence $i\in\lambda$.
        Hence $a\in c(i)$.
        Hence $(\silverIntersectionCode{c}{X}{\lambda})(i)
            \neq(\silverIntersectionCode{c}{Y}{\lambda})(i)$
            by \obv{silver_intersection_codes_differ_at}.
        Suppose $e(X)=e(Y)$.
        $e(X)=\silverIntersectionCode{c}{X}{\lambda}$ and
            $e(Y)=\silverIntersectionCode{c}{Y}{\lambda}$
            by \obv{silver_intersection_map_function}.
        Hence $\silverIntersectionCode{c}{X}{\lambda}
            =\silverIntersectionCode{c}{Y}{\lambda}$.
        Hence $(\silverIntersectionCode{c}{X}{\lambda})(i)
            =(\silverIntersectionCode{c}{Y}{\lambda})(i)$.
        Contradiction.
    \end{subproof}
    Follows by \obv{inj}.
\end{proof}

\begin{theorem}\label{silver_intersection_range_is_eventually_different}
    Let $\lambda$ be a limit ordinal.
    Let $\kappa$ be a limit ordinal.
    Suppose $\normalCofinalMap{c}{\lambda}{\kappa}$.
    Let $A=\silverPowersetFibres{c}{\lambda}$.
    Let $e=\silverIntersectionMap{c}{\kappa}{\lambda}$.
    Let $F=\ran{e}$.
    Then $F\subseteq\dependentProduct{\lambda}{A}$ and
        $\eventuallyDifferentFamily{F}{\lambda}$.
\end{theorem}
\begin{proof}
    $\lambda$ is an ordinal by \obv{limit_ordinal}.
    $c\in\funs{\lambda}{\kappa}$ by \obv{normal_cofinal_map_type}.
    $\cofinalMap{c}{\lambda}{\kappa}$
        by \obv{normal_cofinal_map_is_cofinal}.
    Hence $c$ is a function and $\dom{c}=\lambda$
        by \obv{funs_is_function,funs}.
    Hence $c$ is a function on $\lambda$.
    $e\in\inj{\pow{\kappa}}{\dependentProduct{\lambda}{A}}$
        by \obv{silver_intersection_map_injection}.
    Hence $e\in\funs{\pow{\kappa}}{\dependentProduct{\lambda}{A}}$
        by \obv{inj}.
    Hence $e$ is a function and $\dom{e}=\pow{\kappa}$
        by \obv{funs_is_function,funs}.
    $F\subseteq\dependentProduct{\lambda}{A}$ by \auto{funs_ran}.
    $\ordinalUnbounded{\ran{c}}{\kappa}$
        by \obv{cofinal_map_into_limit_has_unbounded_range}.
    Show that for all $f,g\in F$ such that $f\neq g$ we have
        $\eventuallyDifferent{f}{g}{\lambda}$.
    \begin{subproof}
        Fix $f,g\in F$ such that $f\neq g$.
        Take $X\in\dom{e}$ such that $e(X)=f$
            by \auto{fun_ran_iff}.
        Take $Y\in\dom{e}$ such that $e(Y)=g$
            by \auto{fun_ran_iff}.
        Hence $X\in\pow{\kappa}$ and $Y\in\pow{\kappa}$.
        Hence $X\subseteq\kappa$ and $Y\subseteq\kappa$
            by \obv{powerset_elim}.
        Show that $X\neq Y$.
        \begin{subproof}
            Suppose $X=Y$.
            Hence $e(X)=e(Y)$.
            Hence $f=g$.
            Contradiction.
        \end{subproof}
        Hence $X\symdiff Y$ is inhabited and
            $X\symdiff Y\subseteq\kappa$
            by \obv{distinct_subsets_have_symmetric_difference}.
        Take $a$ such that $a\in X\symdiff Y$.
        Hence $a\in\kappa$ by \obv{subseteq}.
        Take $y\in\ran{c}$ such that $a\in y$
            by \obv{ordinal_unbounded}.
        Take $j\in\dom{c}$ such that $c(j)=y$
            by \obv{fun_ran_iff}.
        Hence $j\in\lambda$.
        Hence $a\in c(j)$.
        $e(X)=\silverIntersectionCode{c}{X}{\lambda}$ and
            $e(Y)=\silverIntersectionCode{c}{Y}{\lambda}$
            by \obv{silver_intersection_map_function}.
        Hence $f=\silverIntersectionCode{c}{X}{\lambda}$ and
            $g=\silverIntersectionCode{c}{Y}{\lambda}$.
        Show that every element of
            $\silverEventualEqualitySet{f}{g}{\lambda}$
            is an element of $\suc{j}$.
        \begin{subproof}
            Fix $i\in\silverEventualEqualitySet{f}{g}{\lambda}$.
            Hence $i\in\lambda$ and $f(i)=g(i)$
                by \obv{silver_eventual_equality_set}.
            Show that $i\in\suc{j}$.
            \begin{subproof}
                Suppose not.
                Hence $i\neq j$ and $i\notin j$ by \auto{suc_iff}.
                $i$ and $j$ are ordinals
                    by \obv{elem_of_ordinal_is_ordinal}.
                Hence $j\in i$ by \obv{ordinal_elem_connex}.
                Hence $c(j)\in c(i)$
                    by \obv{normal_cofinal_map_increasing}.
                $c(i)\in\kappa$ by \obv{funs_type_apply}.
                Hence $c(i)$ is an ordinal
                    by \obv{elem_of_ordinal_is_ordinal,limit_ordinal}.
                Hence $a\in c(i)$ by \obv{ordinal_transitivity}.
                Hence $(\silverIntersectionCode{c}{X}{\lambda})(i)
                    \neq(\silverIntersectionCode{c}{Y}{\lambda})(i)$
                    by \obv{silver_intersection_codes_differ_at}.
                Hence $f(i)\neq g(i)$.
                Contradiction.
            \end{subproof}
        \end{subproof}
        Hence $\silverEventualEqualitySet{f}{g}{\lambda}
            \subseteq\suc{j}$ by \obv{subseteq}.
        $\silverEventualEqualitySet{f}{g}{\lambda}\subseteq\lambda$
            by \obv{silver_eventual_equality_set_subseteq}.
        $\suc{j}\in\lambda$ by \obv{suc_elem_limitordinal}.
        Hence $\ordinalBounded{
            \silverEventualEqualitySet{f}{g}{\lambda}}{\lambda}$
            by \obv{ordinal_bounded}.
        Follows by \obv{eventually_different}.
    \end{subproof}
    Hence $\eventuallyDifferentFamily{F}{\lambda}$
        by \obv{eventually_different_family}.
\end{proof}

\begin{proposition}\label{silver_powerset_fibre_at_gch_point}
    Suppose $c$ is a function on $\lambda$.
    Let $i\in\lambda$.
    Suppose $\gchPoint{c(i)}$.
    Let $A=\silverPowersetFibres{c}{\lambda}$.
    Then $\card{A(i)}=\cardSucc{c(i)}$.
\end{proposition}
\begin{proof}
    $A(i)=\pow{c(i)}$ by \auto{silver_powerset_fibres_type}.
    $\cardPower{\cardTwo}{c(i)}=\cardSucc{c(i)}$
        by \obv{gch_point_elim}.
    $\cardPower{\cardTwo}{c(i)}=\card{\pow{c(i)}}$
        by \obv{cardinal_power_two_is_powerset_cardinal}.
    \begin{align*}
        \card{A(i)}
            &= \card{\pow{c(i)}}
            \\
            &= \cardPower{\cardTwo}{c(i)}
            \\
            &= \cardSucc{c(i)}
    \end{align*}
\end{proof}

\begin{theorem}\label{silver_powerset_upper_bound}
    Let $\kappa$ be a singular cardinal of uncountable cofinality.
    Suppose $\gchPoints{\kappa}$ is stationary in $\kappa$.
    Then $\cardLe{\pow{\kappa}}{\cardSucc{\kappa}}$.
\end{theorem}
\begin{proof}
    Let $\lambda=\cofinality{\kappa}$.
    Hence $\naturals\in\lambda$.
    $\kappa$ is a limit ordinal
        by \obv{infinite_cardinal_is_limit_ordinal}.
    $\cofinality{\kappa}$ is a regular cardinal
        by \obv{cofinality_of_limit_is_regular}.
    Hence $\lambda$ is a regular cardinal.
    Hence $\lambda$ is a limit ordinal
        by \obv{infinite_cardinal_is_limit_ordinal}.
    Hence $\lambda$ is an ordinal by \obv{limit_ordinal}.
    Take $c$ such that
        $\normalCofinalCardinalSequence{c}{\lambda}{\kappa}$ and
        $c$ is index-dominating on $\lambda$ and
        $\infiniteCardinalSequenceValues{c}{\lambda}$
        by \obv{index_dominating_normal_cofinal_cardinal_sequence_exists}.
    $\normalCofinalMap{c}{\lambda}{\kappa}$
        by \obv{normal_cofinal_cardinal_sequence_is_normal_cofinal}.
    Hence $\normalMap{c}{\lambda}{\kappa}$
        by \obv{normal_cofinal_map_is_normal}.
    $\cofinalMap{c}{\lambda}{\kappa}$
        by \obv{normal_cofinal_map_is_cofinal}.
    $\cardinalSequenceValues{c}{\lambda}$
        by \obv{normal_cofinal_cardinal_sequence}.
    $c\in\funs{\lambda}{\kappa}$ by \obv{normal_cofinal_map_type}.
    Hence $c$ is a function and $\dom{c}=\lambda$
        by \obv{funs_is_function,funs}.
    Hence $c$ is a function on $\lambda$.
    Let $A=\silverPowersetFibres{c}{\lambda}$.
    $A$ is a function on $\lambda$
        by \auto{silver_powerset_fibres_type}.
    Let $e=\silverIntersectionMap{c}{\kappa}{\lambda}$.
    Let $F=\ran{e}$.
    $F\subseteq\dependentProduct{\lambda}{A}$ and
        $\eventuallyDifferentFamily{F}{\lambda}$
        by \obv{silver_intersection_range_is_eventually_different}.
    Let $T=\normalMapPreimage{c}{\gchPoints{\kappa}}{\lambda}$.
    $T$ is stationary in $\lambda$
        by \auto{normal_cofinal_map_stationary_pullback}.
    Show that for every $i\in T$ we have
        $\card{A(i)}=\cardSucc{c(i)}$.
    \begin{subproof}
        Fix $i\in T$.
        Hence $i\in\lambda$ and $c(i)\in\gchPoints{\kappa}$
            by \obv{normal_map_preimage}.
        Hence $\gchPoint{c(i)}$ by \obv{gch_points_elim}.
        Follows by \obv{silver_powerset_fibre_at_gch_point}.
    \end{subproof}
    $\powerClosedBelow{\lambda}{\kappa}$
        by \obv{stationary_gch_points_imply_power_closed_below_cofinality}.
    Hence $\cardLe{F}{\cardSucc{\kappa}}$
        by \obv{silver_second_almost_disjoint_bound}.
    $e\in\inj{\pow{\kappa}}{\dependentProduct{\lambda}{A}}$
        by \obv{silver_intersection_map_injection}.
    Hence $e\in\funs{\pow{\kappa}}{\dependentProduct{\lambda}{A}}$
        by \obv{inj}.
    Hence $e\in\surj{\pow{\kappa}}{F}$ by \auto{fun_to_surj}.
    Hence $e\in\funs{\pow{\kappa}}{F}$ by \obv{surj_to_fun}.
    For all $X,Y\in\pow{\kappa}$ such that $e(X)=e(Y)$ we have $X=Y$
        by \obv{inj}.
    Hence $e\in\inj{\pow{\kappa}}{F}$
        by \obv{inj}.
    Hence $\cardLe{\pow{\kappa}}{F}$ by \obv{cardinal_le}.
    Follows by \obv{cardinal_le_transitive}.
\end{proof}

\begin{theorem}[Silver]\label{silver_theorem_stationary}
    Let $\kappa$ be a singular cardinal of uncountable cofinality.
    Suppose $\gchPoints{\kappa}$ is stationary in $\kappa$.
    Then $\cardPower{\cardTwo}{\kappa}=\cardSucc{\kappa}$.
\end{theorem}
\begin{proof}
    $\cardLe{\pow{\kappa}}{\cardSucc{\kappa}}$
        by \obv{silver_powerset_upper_bound}.
    \begin{align*}
        \cardPower{\cardTwo}{\kappa}
            &= \card{\pow{\kappa}}
                \explanation{by \obv{cardinal_power_two_is_powerset_cardinal}}
            \\
            &\subseteq \card{\cardSucc{\kappa}}
                \explanation{by \obv{card_subseteq_iff_cardinal_le}}
            \\
            &= \cardSucc{\kappa}
                \explanation{by \obv{card_of_cardinal,cardinal_successor_is_cardinal}}
    \end{align*}
    Since $\kappa\in\cardPower{\cardTwo}{\kappa}$
        by \obv{cantor_cardinal_power_for_cardinal},
        we have $\cardSucc{\kappa}\subseteq\cardPower{\cardTwo}{\kappa}$
        by \obv{cardinal_element_small,cardinal_power_is_cardinal,cardinal_successor_minimal}.
    Follows by \obv{subseteq_antisymmetric}.
\end{proof}

\begin{corollary}\label{silver_theorem_stationary_powerset}
    Let $\kappa$ be a singular cardinal of uncountable cofinality.
    Suppose $\gchPoints{\kappa}$ is stationary in $\kappa$.
    Then $\card{\pow{\kappa}}=\cardSucc{\kappa}$.
\end{corollary}
\begin{proof}
    $\cardPower{\cardTwo}{\kappa}=\cardSucc{\kappa}$
        by \obv{silver_theorem_stationary}.
    Follows by \auto{cardinal_power_two_is_powerset_cardinal}.
\end{proof}

\begin{corollary}\label{silver_theorem_from_gch_below}
    Let $\kappa$ be a singular cardinal of uncountable cofinality.
    Suppose for every infinite cardinal $\mu\in\kappa$ we have
        $\cardPower{\cardTwo}{\mu}=\cardSucc{\mu}$.
    Then $\cardPower{\cardTwo}{\kappa}=\cardSucc{\kappa}$.
\end{corollary}
\begin{proof}
    Let $\lambda=\cofinality{\kappa}$.
    Hence $\kappa$ is a limit ordinal
        by \obv{infinite_cardinal_is_limit_ordinal}.
    $\lambda$ is a cardinal by \obv{cofinality_is_cardinal}.
    Take $c$ such that
        $\normalCofinalCardinalSequence{c}{\lambda}{\kappa}$ and
        $c$ is index-dominating on $\lambda$ and
        $\infiniteCardinalSequenceValues{c}{\lambda}$
        by \obv{index_dominating_normal_cofinal_cardinal_sequence_exists}.
    $\normalCofinalMap{c}{\lambda}{\kappa}$
        by \obv{normal_cofinal_cardinal_sequence_is_normal_cofinal}.
    $c\in\funs{\lambda}{\kappa}$ by \obv{normal_cofinal_map_type}.
    Hence $c$ is a function and $\dom{c}=\lambda$
        by \obv{funs_is_function,funs}.
    $\clubIn{\ran{c}}{\kappa}$
        by \obv{normal_cofinal_map_range_is_club}.
    Hence $\ran{c}$ is stationary in $\kappa$
        by \obv{club_is_stationary}.
    Show that every element of $\ran{c}$ is an element of
        $\gchPoints{\kappa}$.
    \begin{subproof}
        Fix $\mu\in\ran{c}$.
        Take $i\in\dom{c}$ such that $c(i)=\mu$
            by \obv{fun_ran_iff}.
        Hence $i\in\lambda$.
        Hence $c(i)$ is an infinite cardinal
            by \obv{infinite_cardinal_sequence_value}.
        $\ran{c}\subseteq\kappa$ by \obv{funs_ran}.
        Hence $\mu\in\kappa$ by \obv{subseteq}.
        Hence $\cardPower{\cardTwo}{\mu}=\cardSucc{\mu}$.
        Hence $\gchPoint{\mu}$ by \obv{gch_point}.
        Follows by \obv{gch_points}.
    \end{subproof}
    Hence $\ran{c}\subseteq\gchPoints{\kappa}$ by \obv{subseteq}.
    $\gchPoints{\kappa}\subseteq\kappa$
        by \obv{gch_points_subseteq}.
    Hence $\gchPoints{\kappa}$ is stationary in $\kappa$
        by \obv{stationary_superset}.
    Follows by \obv{silver_theorem_stationary}.
\end{proof}
